\section{Limitations}
\label{sec:limitations}
We discuss several limitations and concerns in this work, revealing potential challenges, constraints, and confusion. 

% First, each challenge in the proposed dataset only consists of 3 positive and 3 negative examples to emphasizing grammar generation in few-shot manner. Typically, other grammar inference algorithms require characteristic samples to uniquely determine the target grammar~\cite{de2010grammatical}. In contrast, we allow general samples that may not uniquely represent a single grammar. Consequently, a given set of examples can correspond to multiple grammars that accept all positives and reject all negatives. This ambiguity motivated the introduction of the $FGI$ metric, which measures how closely a generated grammar approximates the reference grammar. Future work could explore or construct datasets with more examples per challenge or even incorporate characteristic samples to further evaluate LLM performance in grammar generation.

First, although the results indicate that \textit{GPT-4o} exhibits remarkable $SX$ and $SE$, it is important to note that these results may be attributable to the use of GPT-4o during dataset construction. Nonetheless, even though \textit{GPT-4o} already demonstrates excellent performance, \NAME can still enhance it significantly.

Second, as demonstrated, our method does not outperform OPF for \textit{Mistral:7b-Instruct} in $SX$ and $SE$ due to its inherent failure to generate syntactically correct grammars. Nevertheless, our approach yields significant $SX$ and $SE$ improvements for all other LLMs. We also propose to combine syntactical feedback and \NAME to mitigate this limitation and further improve the performance.

Third, one may argue that given any finite set of positive and negative examples, it could always be possible to construct a regular grammar rather than a CFG that can accept all positives and reject all negatives. However, such an approach may function more like a classifier rather than a grammar and may lack applicability in subsequent tasks, such as constructing an abstract syntax tree. 

Finally, in this work, we primarily focus on LLM-based few-shot grammar generation without comparing algorithms that are not LLM-based. The reasons behind this are that most algorithms require a large set of characteristic examples to uniquely determine the target grammar~\cite{de2010grammatical}. Instead, we do not impose such constraints on our example set and hypothesize that the experience and knowledge acquired from corpus can enable LLMs to handle few-shot grammar generation tasks. Consequently, those algorithms may not be directly applicable. In addition, since we focus on the exploration and improvement of the ability of LLMs in few-shot grammar generation, we construct two LLM-based baselines for fair comparison. 
